\section{OWASP-Jucieshop Code-analyse \label{cha:demo}} 
Wie in Kapitel \ref{sub:Code_Analyse} beschrieben, bezeichnet die Schwachstellenanalyse die systematische Untersuchung eines digitalen Systems auf sicherheitsrelevante Fehler. Bei einer quellcodebasierten Analyse wird die Implementierung einer Anwendung untersucht, um potenzielle Schwachstellen und deren Ursachen zu identifizieren. In diesem Kapitel wird anhand eines praktischen Versuchs untersucht, inwieweit ein generatives KI Modell eine solche Schwachstelle erkennen, hinsichtlich ihrer Ursache einordnen und einen möglichen Angriffsweg erläutern kann. Als Testumgebung dient der OWASP Juice Shop, eine absichtlich verwundbare Webanwendung, die für Schulungs- und Forschungszwecke entwickelt wurde.
 
Die Testumgebung besteht aus einer lokal betriebenen Instanz des OWASP Juice Shops, die innerhalb eines Docker Containers ausgeführt wird. Für die Analyse wurde das generative KI Modell Claude Opus 4.6 eingesetzt. Die Untersuchung konzentriert sich auf die Authentifizierungslogik der Anwendung. Ziel ist es, zu prüfen, ob das Modell eine im Quellcode enthaltene Schwachstelle identifizieren, die betroffenen Dateien und Funktionen lokalisieren sowie die technische Ursache und mögliche Gegenmaßnahmen erläutern kann. Dem Modell wurde dazu der Quellcode der Anwendung in einer kontrollierten, ausschließlich lokalen Testumgebung zur Verfügung gestellt. Der verwendete Prompt lautete wie folgt:

\begin{quote}
    \textit{You are assisting in an authorized security training lab using the given webshop source code located in \texttt{./project}. Review only the provided codebase (no scanning, brute forcing, or interaction with live systems). Locate the authentication/login implementation, trace how username/email and password are processed, identify the vulnerability enabling the first login-bypass challenge, and cite the relevant file(s), function(s), and vulnerable code pattern. Explain the root cause, classify it using OWASP terminology, and describe why the implementation is insecure. Then provide progressive hints for solving the challenge without immediately revealing the payload, followed by a high-level explanation of how the bypass works and a secure remediation using parameterized queries, prepared statements, or safe ORM practices. If the repository is unavailable, specify exactly which files are needed.}
\end{quote}

Nach der Einrichtung der Testumgebung wurde der Prompt an das Modell übermittelt. Claude Opus 4.6 identifizierte die Schwachstelle in der Authentifizierungslogik und verwies auf die betroffenen Dateien und Funktionen. Das Modell erkannte, dass Benutzereingaben innerhalb der Login Funktion auf unsichere Weise verarbeitet wurden, wodurch ein Injection Angriff ermöglicht wurde. Darüber hinaus erläuterte es die technische Ursache der Schwachstelle, ordnete sie einer entsprechenden OWASP Kategorie zu und beschrieb geeignete Gegenmaßnahmen.

Die Hinweise des Modells ermöglichten es, die Schwachstelle innerhalb der lokalen Testumgebung nachzuvollziehen und die Authentifizierung als administrativer Benutzer zu umgehen. Ursache hierfür war die unsichere Verarbeitung der Anmeldedaten. Als Gegenmaßnahme schlug das Modell unter anderem die Verwendung parametrisierter Abfragen, vorbereiteter Anweisungen und sicherer ORM Funktionen vor. Damit konnte das Modell nicht nur die Schwachstelle identifizieren, sondern auch geeignete Maßnahmen zu ihrer Behebung nennen.

Das Ergebnis des Versuchs zeigt, dass das untersuchte generative KI Modell in diesem konkreten Szenario in der Lage war, eine Schwachstelle anhand des bereitgestellten Quellcodes zu identifizieren, ihre Ursache zu erläutern und geeignete Gegenmaßnahmen vorzuschlagen. Gleichzeitig demonstriert der Versuch den Dual Use Charakter solcher Systeme: Dieselben Fähigkeiten können sowohl zur Unterstützung legitimer Sicherheitsanalysen als auch zur Vorbereitung missbräuchlicher Angriffe eingesetzt werden. Aufgrund des begrenzten Versuchsumfangs lassen sich aus diesem Einzelergebnis jedoch keine allgemeingültigen Aussagen über die Zuverlässigkeit generativer KI Modelle bei der Schwachstellenanalyse ableiten.
% \section{Fazit}
% Durch den einsatz von KI Hilfsmittelen von Angreifern ist es leichter Schwachstellen zu finden und diese auszunutzen. Gerade im Bereich Phishing und Socialengineering macht der gebrauch einer LLM vieles einfacher und schneller. Um eine gute qualität zu erhalten ist kein Top Modell erfoderlich, wie z.B. bei der Schwachstellenanalyse. Da Software immer sicherer wird und die meisten Angrife durch Phishing oder Socialengineering gestartet werden, senkt ein LLM diese Hürde noch weiter. Trotzdem ändert ein LLM nicht unbedingt den Ablauf sondern vereinfach nur Schritte, wie das übersetzen. Die Gefahr steigt für Privat Personen mehr and als für Firmen. Da im Privatenbereich nicht so viele Schutzmechanismen existieren wie im Enterprisebereich.

% Die interessanteren Angriffsvektoren sind die Schwachstellenanalyse von . 
% Unter betrachtung der Ergebnisse aus Kapitel \ref{cha:mythosvsoffentlich} und dem eigenem Versuch im Kapitel \ref{cha:demo}. Daraus kann abgeleitet werden das die Ergebnisse von Mythos eine bessere Qualität haben. Die Schwachstellen könnnen auch mit öffentlich zugänglichen Modellen gefunden und ausgenutzt werden, die Qualität und konsetenz ist nicht so sicher wie bei Mythos.
